% ------------------------------------------------------------------------
% Modelo de Trabalho de Conclusão de Curso em conformidade com 
% ABNT NBR 14724:2011: Informacao e documentacao - Trabalhos academicos -
% Apresentacao:
% Versão: 2026-07-30
% Atualizado por: frozza@gmail.com
% ------------------------------------------------------------------------

\documentclass[12pt, oneside, a4paper, brazil]{abntex2}
% ---
% Pacotes básicos 
% ---
\usepackage{helvet}			     % Usa a fonte Arial-like	
\renewcommand{\familydefault}{\sfdefault} % Usa a fonte Arial-like
\usepackage[T1]{fontenc}		 % Selecao de codigos de fonte.
\usepackage[utf8]{inputenc}		 % Codificacao do documento (conversão automática dos acentos)
\usepackage{lastpage}			 % Usado pela Ficha catalográfica
\usepackage{indentfirst}		 % Indenta o primeiro parágrafo de cada seção.
\usepackage{xcolor}			     % Controle das cores
\usepackage{graphicx}			 % Inclusão de gráficos
\usepackage{microtype} 			 % para melhorias de justificação
\usepackage[alf]{abntex2cite}	 % Citações padrão ABNT
\usepackage{lipsum}              % Pode ser removido no final

%===== PACKAGES FROM AUTHORS
\newcommand{\hl}[1]{\colorbox{yellow}{#1}} % para realçar o texto
\usepackage{comment}
\usepackage{tabularx}            % Necessário para o ambiente tabularx
\usepackage{array}               % Geralmente incluído com tabularx, mas bom para p, m, b
\usepackage{float}
\usepackage{pdflscape}           % para inserir páginas com orientação paisagem
\usepackage[section]{placeins}   % não deixa figura atravessar sections
% hyperref deve ficar sempre por último
\usepackage{hyperref}            % Para o comando \url

%%%% Configuracao de Quadros
%===== PARA USAR QUADROS (versão ajustada e mais robusta) =====
\usepackage{float}      % para criar novos ambientes flutuantes
\usepackage{caption}    % para configurar o caption

% ===== AMBIENTE QUADRO =====

\newfloat{quadro}{htbp}{loq}
\floatname{quadro}{Quadro}

% Numeração simples
\renewcommand{\thequadro}{\arabic{quadro}}

% Lista de Quadros
\providecommand{\listquadroname}{Lista de Quadros}
\newcommand{\listofquadro}{\listof{quadro}{\listquadroname}}

% Formato do caption
\DeclareCaptionLabelFormat{quadronum}{#1~#2\ - }

\captionsetup[quadro]{
    name=Quadro,
    labelformat=quadronum,
    labelsep=none,
    position=top,
    font=normal,
    labelfont=bf
}


%%%%% Configuração de Listas de código
\usepackage{listings}
\usepackage{xcolor}


% ----------------------------------------------------------
% CORES
% ----------------------------------------------------------
\definecolor{codeblue}{RGB}{33,74,135}
\definecolor{codegreen}{RGB}{0,120,0}
\definecolor{codegray}{RGB}{100,100,100}
\definecolor{codeorange}{RGB}{180,90,0}
\definecolor{codebg}{RGB}{245,245,245}

% ----------------------------------------------------------
% SUPORTE A UTF-8 (ACENTOS)
% ----------------------------------------------------------
\lstset{
  inputencoding=utf8,
  extendedchars=true,
  literate=
    {á}{{\'a}}1 {à}{{\`a}}1 {ã}{{\~a}}1 {â}{{\^a}}1
    {é}{{\'e}}1 {ê}{{\^e}}1
    {í}{{\'i}}1
    {ó}{{\'o}}1 {ô}{{\^o}}1 {õ}{{\~o}}1
    {ú}{{\'u}}1
    {ç}{{\c c}}1
    {Á}{{\'A}}1 {Ã}{{\~A}}1 {Â}{{\^A}}1
    {É}{{\'E}}1 {Ê}{{\^E}}1
    {Í}{{\'I}}1
    {Ó}{{\'O}}1 {Ô}{{\^O}}1 {Õ}{{\~O}}1
    {Ú}{{\'U}}1
    {Ç}{{\c C}}1
}

% ----------------------------------------------------------
% ESTILO BASE (ROBUSTO)
% ----------------------------------------------------------
\lstdefinestyle{base}{
    basicstyle=\ttfamily\small,
    backgroundcolor=\color{codebg},
    breaklines=true,
    breakatwhitespace=false,
    columns=fullflexible,
    keepspaces=true,
    frame=single,
    numbers=left,
    numberstyle=\tiny\color{codegray},
    numbersep=6pt,
    showstringspaces=false,
    tabsize=2,
    captionpos=t,
    aboveskip=1em,
    belowskip=1em,
    xleftmargin=0.5em,
    xrightmargin=0.5em,
    postbreak=\mbox{\textcolor{codegray}{$\hookrightarrow$}\space}
}

% ==========================================================
% JAVASCRIPT
% ==========================================================
\lstdefinelanguage{JavaScript}{
  keywords={
    break, case, catch, continue, debugger, default, delete, do, else,
    finally, for, function, if, in, instanceof, new, return, switch,
    this, throw, try, typeof, var, void, while, with, const, let,
    async, await
  },
  keywordstyle=\color{codeblue}\bfseries,
  comment=[l]{//},
  morecomment=[s]{/*}{*/},
  commentstyle=\color{codegreen}\itshape,
  stringstyle=\color{codeorange},
  morestring=[b]',
  morestring=[b]"
}

\lstdefinestyle{javascript}{
    style=base,
    language=JavaScript
}

% ==========================================================
% JSON (melhorado)
% ==========================================================
\lstdefinestyle{json}{
    style=base,
    showstringspaces=false,
    basicstyle=\ttfamily\small
}

% ==========================================================
% BASH / SHELL
% ==========================================================
\lstdefinestyle{bashblock}{
    style=base,
    language=sh
}

% ==========================================================
% JAVA
% ==========================================================
\lstdefinestyle{java}{
    style=base,
    language=Java,
    keywordstyle=\color{codeblue}\bfseries,
    commentstyle=\color{codegreen}\itshape,
    stringstyle=\color{codeorange}
}

% ==========================================================
% PYTHON
% ==========================================================
\lstdefinestyle{python}{
    style=base,
    language=Python,
    keywordstyle=\color{codeblue}\bfseries,
    commentstyle=\color{codegreen}\itshape,
    stringstyle=\color{codeorange},
    morekeywords={self, cls}
}

% ==========================================================
% C++
% ==========================================================
\lstdefinestyle{cpp}{
    style=base,
    language=C++,
    keywordstyle=\color{codeblue}\bfseries,
    commentstyle=\color{codegreen}\itshape,
    stringstyle=\color{codeorange}
}

% ==========================================================
% DOCKER
% ==========================================================
\lstdefinestyle{dockerblock}{
    style=base,
    language=sh
}

% ==========================================================
% SQL
% ==========================================================
\lstdefinestyle{sql}{
    style=base,
    language=SQL,
    keywordstyle=\color{codeblue}\bfseries,
    commentstyle=\color{codegreen}\itshape,
    stringstyle=\color{codeorange},
    morekeywords={
        PRIMARY, KEY, NOT, NULL, UNIQUE, DEFAULT,
        DATE, SERIAL, CREATE, TABLE, INSERT,
        UPDATE, DELETE, SELECT, FROM, WHERE
    }
}

% ==========================================================
% XML
% ==========================================================

\lstdefinelanguage{XML}{
    morestring=[b]",
    morestring=[b]',
    stringstyle=\color{codeorange},
    comment=[s]{<!--}{-->},
    commentstyle=\color{codegreen}\itshape,
    morekeywords={<?xml, xmlns, version, encoding},
    keywordstyle=\color{codeblue},
    alsoletter={-},
    sensitive=true
}

\lstdefinestyle{xml}{
    style=base,
    language=XML
}

% ==========================================================
% ABNT
% ==========================================================
\renewcommand{\lstlistingname}{Código}
\renewcommand{\lstlistlistingname}{Lista de códigos}

% Padroniza com abnTeX2
\makeatletter
\renewcommand{\lstlistoflistings}{%
  \begingroup
  \phantompart
  \pdfbookmark[0]{\lstlistlistingname}{lol}
  \listof{lstlisting}{\lstlistlistingname}
  \endgroup
}
\makeatother

% Numeração contínua
\lstset{numberbychapter=false}

% Lista personalizada com "Código X -"
\newcommand{\listofcodigos}{
  \begingroup
    \renewcommand{\numberline}[1]{Código~##1~--~}
    \lstlistoflistings
  \endgroup
}

\titulo{Tema do TCC}
\autor{Aluno 1\\ Aluno 2\\ Outro aluno}

\local{Campina Grande -- PB}
\data{dezembro de 2026}
\orientador{Prof. Dr. fulano de tal}
%\coorientador{Prof. Dr. Outro Fulano Coorientador}

\makeatletter
\renewcommand{\imprimircapa}{%
  \begin{capa}
    \center

    % ================= LOGO =================
    %\vspace*{1cm}
    \includegraphics[width=3cm]{Figuras/Instituto_Federal_da_Paraíba_-_Logo.png}

    %\vspace{1cm}

    % ================= CABEÇALHO =================
    {\ABNTEXchapterfont\small
    Instituto Federal da Paraíba\\
    Bacharelado em Engenharia de Computação\\
    Campus Campina Grande
    \par}

    \vfill

    % ================= AUTOR =================
    {\ABNTEXchapterfont\large\MakeUppercase{\imprimirautor}\par}

    \vfill

    % ================= TÍTULO =================
    {\ABNTEXchapterfont\bfseries\Large\MakeUppercase{\imprimirtitulo}\par}

    \vfill

    % ================= LOCAL E ANO =================
    {\large
      \imprimirlocal\\
      \imprimirdata
      \par}

  \end{capa}
}
\makeatother
          % Configuração da capa
% ------------------------------------------------------------------------
% Folha de Rosto personalizada (Autor e Título em MAIÚSCULAS)
% Substitui o ambiente folhaderosto do abnTeX2
% \begin{folhaderosto}
% \end{folhaderosto}
% ------------------------------------------------------------------------

\makeatletter

% ---- Ambiente principal ----
\renewenvironment{folhaderosto}
{%
%  \cleardoublepage
\clearpage

  \thispagestyle{empty}

  \begin{center}

  % Autor
  {\ABNTEXchapterfont\large\MakeUppercase{\imprimirautor}\par}

  \vspace*{\fill}

  % Título
  {\ABNTEXchapterfont\bfseries\Large\MakeUppercase{\imprimirtitulo}\par}

  \vspace*{\fill}

  % Preâmbulo + orientações
  \hspace{.45\textwidth}
  \begin{minipage}{.5\textwidth}
    \imprimirpreambulo\par

    \vspace{1em}

    \imprimirorientadorRotulo~\imprimirorientador

%    \ifnotempty{\imprimircoorientador}{%
%      \par
%      \imprimircoorientadorRotulo~\imprimircoorientador
%    }

      \ifx\imprimircoorientador\empty
          % não faz nada
      \else
         \par
         \imprimircoorientadorRotulo~\imprimircoorientador
      \fi
  \end{minipage}

  \vspace*{\fill}

  % Local e data
  {\large\imprimirlocal\par}
  {\large\imprimirdata\par}

  \end{center}

%  \cleardoublepage
}
{}
\makeatother
  % Configuração Folha de Rosto

% Configurações de aparência do PDF final
% informações do PDF
\makeatletter
\hypersetup{
		pdftitle={\@title}, 
		pdfauthor={\@author},
   	  pdfsubject={\imprimirpreambulo},
        pdfcreator={LaTeX with abnTeX2},
		pdfkeywords={abnt}{latex}{abntex}{abntex2}{trabalho acadêmico}, 
		colorlinks=true,       		% false: boxed links; true: colored links
     	linkcolor=blue,          	% color of internal links
    	citecolor=blue,        		% color of links to bibliography
    	filecolor=magenta,     		% color of file links
		urlcolor=blue,
		bookmarksdepth=4
}
\makeatother

% Espaçamentos entre linhas e parágrafos 
\setlength{\parindent}{1.3cm}			% O tamanho do parágrafo é dado por:
\setlength{\parskip}{0cm}  			% Controle do espaçamento entre um parágrafo e outro:


% Início do documento
\begin{document}
    \selectlanguage{brazil} 				% Seleciona o idioma do documento (conforme pacotes do babel)
    \frenchspacing 							    % Retira espaço extra obsoleto entre as frases.
    
    % ELEMENTOS PRÉ-TEXTUAIS
    \imprimircapa							% Capa

    % Inserir Folha de Rosto
    \begin{folhaderosto}
    \end{folhaderosto}    
    
    % Inserir a ficha bibliografica
    
    % Isto é um exemplo de Ficha Catalográfica, ou ``Dados internacionais de 
    % catalogação-na-publicação''. Utilizar este modelo como referência. 
    % Porém, provavelmente a biblioteca fornecerá um PDF com a ficha catalográfica 
    % definitiva após a defesa do trabalho. Quando estiver com o documento, salve-o como PDF 
    % no diretório do seu projeto e substitua todo o conteúdo de implementação deste arquivo 
    % pelo comando abaixo:
    %
    % \begin{fichacatalografica}
    %     \includepdf{fig_ficha_catalografica.pdf}
    % \end{fichacatalografica}
    \clearpage
    \input{Pre-Texto/fichacatalografica}
    
    
    % Inserir folha de aprovação
    
    % Exemplo de Folha de aprovação, elemento obrigatório da NBR 14724/2011 (seção 4.2.1.3). 
    % Utilizar este modelo até a aprovação do trabalho. Após isso, substitua todo o conteúdo 
    % deste arquivo por uma imagem da página assinada pela banca com o comando abaixo:
    %
    % \includepdf{folhadeaprovacao_final.pdf}
    %
    % Inserir folha de aprovação

% Exemplo de Folha de aprovação, elemento obrigatório da NBR 14724/2011 (seção 4.2.1.3). 
% Utilizar este modelo até a aprovação do trabalho. Após isso, substitua todo o conteúdo 
% deste arquivo por uma imagem da página assinada pela banca com o comando abaixo:
%
% \includepdf{folhadeaprovacao_final.pdf}
%
\begin{folhadeaprovacao}

  \begin{center}
    {\ABNTEXchapterfont\large\imprimirautor}

    \vspace*{\fill}\vspace*{\fill}
    \begin{center}
      \ABNTEXchapterfont\bfseries\Large\imprimirtitulo
    \end{center}
    \vspace*{\fill}
    
    \hspace{.45\textwidth}
    \begin{minipage}{.5\textwidth}
        \imprimirpreambulo
    \end{minipage}%
    \vspace*{\fill}
   \end{center}
        
   Trabalho aprovado. \imprimirlocal, 17 de dezembro de 2026:

   \assinatura{\textbf{\imprimirorientador} \\ Orientador} 
   \assinatura{\textbf{Maria Fulana, Dra.} \\ Membro banca 1}
   %\assinatura{\textbf{Fulaninho Silva, Dr.} \\ Membro banca 2}
   %\assinatura{\textbf{Professor} \\ Convidado 3}
   %\assinatura{\textbf{Professor} \\ Convidado 4}
      
   \begin{center}
    \vspace*{0.5cm}
    {\large\imprimirlocal}
    \par
    {\large\imprimirdata}
    \vspace*{1cm}
  \end{center}
  
\end{folhadeaprovacao}
% ---

%% Folha de assinaturas
Assinaturas digitais
\newpage
  
    \input{Pre-Texto/dedicatoria}
    \input{Pre-Texto/resumo}
    \input{Pre-Texto/indices-siglas}
    
    % ----------------------------------------------------------
    % ELEMENTOS TEXTUAIS
    % ----------------------------------------------------------
    \textual
    
    % ---
    % Inclusão de capítulo de Introdução ao trabalho
    % ---
    \input{Capitulos/1-Introducao}
    % ---
    % Inclusão de outros capitulos
    % ---
    \input{Capitulos/2-Referencial}
    \input{Capitulos/3-AnaliseEProjeto}
    \input{Capitulos/4-Implementacao}
    
    % ----------------------------------------------------------
    % Finaliza a parte no bookmark do PDF
    % para que se inicie o bookmark na raiz
    % e adiciona espaço de parte no Sumário
    % ----------------------------------------------------------
    \phantompart
    
    \input{Capitulos/5-Conclusao}
    
    % ----------------------------------------------------------
    % ELEMENTOS PÓS-TEXTUAIS
    % ----------------------------------------------------------
    \postextual
    % ----------------------------------------------------------
    
    % ----------------------------------------------------------
    % Referências bibliográficas
    % ----------------------------------------------------------
    \bibliography{referencias}
    
    % ----------------------------------------------------------
    % Glossário
    % ----------------------------------------------------------
    %
    % Consulte o manual da classe abntex2 para orientações sobre o glossário.
    %
    %\glossary
    
    % ----------------------------------------------------------
    % Apêndices
    % ----------------------------------------------------------
    
    % ---
    % Inicia os apêndices
    % ---
    \begin{apendicesenv}
    
    % Imprime uma página indicando o início dos apêndices
    \partapendices
    
    % ----------------------------------------------------------
    \chapter{Quisque libero justo}
    % ----------------------------------------------------------
    
    \lipsum[50]
    
    % ----------------------------------------------------------
    \chapter{Nullam elementum urna vel imperdiet sodales elit ipsum pharetra ligula
    ac pretium ante justo a nulla curabitur tristique arcu eu metus}
    % ----------------------------------------------------------
    \lipsum[55-57]
    
    \end{apendicesenv}
    % ---
    
    
    % ----------------------------------------------------------
    % Anexos
    % ----------------------------------------------------------
    
    % ---
    % Inicia os anexos
    % ---
    \begin{anexosenv}
    
    % Imprime uma página indicando o início dos anexos
    \partanexos
    
    % ---
    \chapter{Morbi ultrices rutrum lorem.}
    % ---
    \lipsum[30]
    
    % ---
    \chapter{Cras non urna sed feugiat cum sociis natoque penatibus et magnis dis
    parturient montes nascetur ridiculus mus}
    % ---
    
    \lipsum[31]
    
    % ---
    \chapter{Fusce facilisis lacinia dui}
    % ---
    
    \lipsum[32]
    
    \end{anexosenv}
\end{document}
